%------------------------------------------------------------------------------%
\begin{question}
  Represente matricialmente a equação
  \[
    x^2 + y^2 = 9
  \]
\end{question}

%------------------------------------------------------------------------------%
\begin{resolution}{Solução}
\begin{columns}[t]
\column{0.5\textwidth}
  Equação
  \[
    x^2 + y^2 = 9
  \]

  \pause
  Parte quadrática
  \[
    \pause
    ax^2 + 2bxy + cy^2
    \pause
    = x^2 + 2y^2
  \]
  \[
    \pause
    A = \begin{bmatrix}
      1 & 0 \\
      0 & 2
    \end{bmatrix}
  \]

\column{0.4\textwidth}
  \pause
  Termos lineares
  \[
    \pause
    p = \begin{bmatrix} 0 \\ 0 \end{bmatrix}
  \]

  \bigskip
  \pause
  Termo constante
  \[
    f = -9
  \]
\end{columns}
\end{resolution}

%------------------------------------------------------------------------------%
\endinput

\begin{frame}{Exemplo 3 — Enunciado}
Represente matricialmente a equação
\[
x^2+y^2=9
\]
\end{frame}

\begin{frame}{Exemplo 3 — Resolução}
Definimos
\[
A=\begin{bmatrix}1&0\\0&1\end{bmatrix}
\qquad
\vec{x}=\begin{bmatrix}x\\y\end{bmatrix}
\]
Então
\[
\vec{x}^{\,t}A\vec{x}=9
\]
pois
\[
\begin{bmatrix}x&y\end{bmatrix}
\begin{bmatrix}1&0\\0&1\end{bmatrix}
\begin{bmatrix}x\\y\end{bmatrix}
=x^2+y^2
\]
A equação é a de uma circunferência centrada na origem.
\end{frame}
